% week10_thu_show_and_tell.tex -- DSCT Week 10 Thursday, Oct 22 2026.
% Sources: week-10/instructor/thursday-run-of-show.md, week-10/README.md,
% assignments/show-and-tell-artifact.md (critique norm), docs/grading-model.md.
% Numbers on the warm-up slide were recomputed in Python.
% NOTE for the instructor: run-of-show and README still say "draft/submit the
% individual Show & Tell report", but Decision 029 (2026-08-25) retired that
% report. This deck follows Decision 029: Show & Tell is in-class and ungraded;
% the Decision Gate is the one submission. Edit slide 13 if that is wrong.
\documentclass[aspectratio=169,10pt]{beamer}
\usetheme{Madrid}
\title{Defend a Probability Claim}
\subtitle{Week 10 Thursday --- Show \& Tell}
\author{COMSC 2043 \textbullet\ Discrete Structures and Critical Thinking}
\date{October 22, 2026}
\begin{document}

\frame{\titlepage
\note{75 minutes. Protect critique and revision time. A hand calculation is sufficient; reward inspectable evidence, not polished slides. Run of show: 0:00 denominator check swap (10). 0:10 defense and critique norms (8). 0:18 rotating defenses (30). 0:48 peer questions and professional critique (15). 1:03 revise conclusion and boundary (13). 1:16 separate Show and Tell from the Decision Gate, preview relations (9).}}

\begin{frame}{Today's destination}
\begin{itemize}
\item Swap one denominator check.
\item Defend a claim in four beats: claim, table or simulation, assumption, boundary.
\item Ask and receive professional critique of the reasoning, not the reasoner.
\item Revise your conclusion so it is no stronger than the evidence.
\end{itemize}
\note{Show and Tell is public explanation and defense under questions, not presentation theater. Revision under challenge is not failure.}
\end{frame}

\begin{frame}{Warm-up: check a classmate's denominator}
\begin{block}{Scenario}
1{,}000 people. Prevalence 2\%. Sensitivity 80\%. Specificity 90\%.\\
A classmate reports: ``$P(C\mid +) = 16/20 = 80\%$.''
\end{block}
\begin{itemize}
\item Swap with a partner. Is the denominator right for $P(C\mid +)$?
\item Rebuild the four-cell table, then confirm or correct.
\end{itemize}
\note{Ten minutes. Verified in Python: C = 20 people, 16 test positive, 4 negative; not-C = 980, of whom 98 test positive and 882 negative; cells total 1,000; positives = 114. Correct answer is 16/114 = 8/57, about 14.0 percent. The classmate's 16/20 = 80 percent is sensitivity P(+ given C), the wrong denominator. Corrected or confirmed check is the visible work.}
\end{frame}

\begin{frame}{The correct check}
\centering
\begin{tabular}{l|r|r|r}
 & Tests $+$ & Tests $-$ & Total \\ \hline
Has $C$ (2\%) & 16 & 4 & 20 \\
Not $C$ (98\%) & 98 & 882 & 980 \\ \hline
Total & 114 & 886 & 1{,}000
\end{tabular}
\vspace{3mm}
\begin{itemize}
\item $P(C\mid +)$: denominator is the $+$ column, 114. So $16/114=8/57\approx 14.0\%$.
\item $16/20=80\%$ answers a different question: $P(+\mid C)$.
\end{itemize}
\note{Cells: 16 + 4 + 98 + 882 = 1,000. Negatives column total is 4 + 882 = 886. Make sure every pair says which column is the denominator out loud.}
\end{frame}

\begin{frame}{What Show \& Tell is}
\begin{itemize}
\item You make a claim public and defend it under questions.
\item Not presentation theater: a hand calculation is enough.
\item The evidence is the point. Polished slides earn nothing extra.
\item You also give at least one specific, professional critique to a peer.
\end{itemize}
\note{Pull from the Show and Tell contract: claim, strongest evidence, strongest challenge, what survived or changed, evidence boundary, critique given. Reward inspectable evidence.}
\end{frame}

\begin{frame}{The critique norm}
\begin{block}{Critique the claim, assumption, evidence, inference, or model.}
Do not attack the reasoner.
\end{block}
\begin{itemize}
\item Agreeing is not an exemption: probe whether the evidence deserves the confidence claimed.
\item Revising a claim under challenge is not failure.
\item Refusing to revise when the evidence demands it is not strength.
\end{itemize}
\note{Read the norm aloud. Model one good critique and one bad one (the bad one is about the person). Do not reward disagreement for its own sake or rhetorical dominance.}
\end{frame}

\begin{frame}{Questions that expose something}
\begin{itemize}
\item What is the denominator, and why that one?
\item Which assumption, if false, breaks this number? (representative population, independence, randomness)
\item What did the simulation or second calculation check, and what could it never check?
\item Is the conclusion stronger than the table supports?
\item Where did the AI help, and how did you check it independently?
\end{itemize}
\note{A starter list, not a script. Encourage students to say why the question matters: what uncertainty, risk, or missing evidence it is meant to expose.}
\end{frame}

\begin{frame}{Your defense: four beats}
\begin{enumerate}
\item \textbf{Claim} --- one sentence.
\item \textbf{Table or simulation} --- show the counts or the code and the check.
\item \textbf{Assumption} --- name the one that matters most.
\item \textbf{Boundary} --- what your evidence does not justify.
\end{enumerate}
\vspace{2mm}
\small Then take questions. Keep the artifact inspectable: a hand-drawn table is fine.
\note{Rotating defenses run 0:18 to 0:48. Keep time by signal, not by clock-watching: when the signal sounds, the listener becomes the defender or the group rotates. Do not let a gambling example become an unstated real-world frequency claim.}
\end{frame}

\begin{frame}{Rotating defenses}
\Large\centering
\textbf{Defend $\rightarrow$ question $\rightarrow$ rotate.}\\[4mm]
\normalsize
Take notes on the strongest challenge you receive and the critique you give.
\vspace{2mm}
\begin{block}{Capture as you go}
Strongest challenge received. What survived or changed. One critique you gave and why it was useful.
\end{block}
\note{Thirty minutes. Circulate, listen for denominators and independence assumptions, and note one or two for the whole-class debrief. These notes are working notes for revision.}
\end{frame}

\begin{frame}{Peer questions and professional critique}
\begin{itemize}
\item Ask the question you would want asked about your own claim.
\item Name the thing you are challenging: claim, assumption, evidence, inference, or model.
\item Say why it matters, then listen to the answer.
\end{itemize}
\vspace{2mm}
\begin{alertblock}{Not the same as an attack}
``Your 90\% is sensitivity, not $P(C\mid +)$'' targets the inference. ``You never get this stuff'' targets the person.
\end{alertblock}
\note{0:48 to 1:03. Fifteen minutes of whole-class or large-group questions. Bring up the strongest critiques you overheard. Keep it about reasoning.}
\end{frame}

\begin{frame}{Revise: make the conclusion no stronger than the evidence}
\begin{columns}[T]
\column{.48\textwidth}
\textbf{Before}\\[1mm]
``A positive result means you probably have the condition.''
\column{.48\textwidth}
\textbf{After}\\[1mm]
``In the 10{,}000-person example with 1\% prevalence, about 15 of every 100 positives have the condition ($90/585$). This holds only if the stated rates apply to the population.''
\end{columns}
\vspace{3mm}
\small The revision names the base rate, the number, and the boundary.
\note{Numbers re-verified: 90/585 = 15.4 percent at 1 percent prevalence from Tuesday's table. Students now revise their own conclusion and boundary using what they heard.}
\end{frame}

\begin{frame}{Work time: revise your conclusion and boundary}
\begin{itemize}
\item Rewrite your one-sentence conclusion after the challenges you heard.
\item State what survived, what you narrowed, and what remains unresolved.
\item Note where AI helped and how you independently checked it.
\end{itemize}
\vspace{2mm}
\begin{block}{Where this goes}
These revisions feed your Decision Gate: Sources, Rules/Assumptions, Work, Check, One-Sentence Summary.
\end{block}
\note{1:03 to 1:16. Protect this time; do not let the debrief eat it. Walk around and ask each student which assumption they narrowed.}
\end{frame}

\begin{frame}{Show \& Tell and the Decision Gate are different things}
\begin{itemize}
\item \textbf{Show \& Tell} is today's in-class practice: defend, question, critique, revise. No separate Canvas assignment.
\item \textbf{Decision Gate} is the one submission this week: an individual write-up with the five headings.
\item Do not submit the same artifact twice, and do not write an extra essay.
\end{itemize}
\small The due date and time are on the Decision Gate in Canvas.
\note{INSTRUCTOR CHECK: the Week 10 README and run-of-show still say draft or submit an individual Show and Tell report, but Decision 029 (2026-08-25) retired that report and docs/grading-model.md calls Show and Tell ungraded and in-class. This slide follows Decision 029. If you still want a short written report this week, change the first bullet and tell the class where to submit it.}
\end{frame}

\begin{frame}{Before you leave: your submission plan}
\begin{enumerate}
\item Which claim will your Decision Gate defend? (screening example, a stated game or drawing process, or an approved bounded claim)
\item What is your Check, and what can it not validate?
\item Which of today's critiques changed your Rules/Assumptions?
\item When will you draft the one-sentence summary?
\end{enumerate}
\note{1:16 to 1:25. Collect a one-line submission plan from each student or have them hold it. The plan is retained, not graded.}
\end{frame}

\begin{frame}{Next week: relations}
\begin{itemize}
\item From probability over events to relations over pairs.
\item Which pairs are related, and how do we classify the relation?
\item Keep your Week 10 trace: it feeds the Decision Gate.
\end{itemize}
\note{Bridge to Week 11: relations, equivalence, partial orders, matrices and digraphs. Keep it to a preview.}
\end{frame}

\begin{frame}{Exit question}
\Large\centering
\textbf{What is one claim you narrowed today, and what did the evidence not justify?}
\note{Collect answers verbally or on paper as the room clears.}
\end{frame}

\end{document}
